\section{Stability and containment}
\label{sec:stability}

A credit community fails in one of two ways: it underwrites more than it can absorb, or
one member's failure propagates until many fail. The mechanism of \S\ref{sec:standing}
does most of the work against both.

\subsection{Aggregate exposure is bounded by construction}

The total insured credit a community can extend is not a quantity anybody sets. By
Corollary~\ref{cor:uw-bound} it is bounded by the sum of the supplies its underwriters
have declared, a figure the community explicitly committed to. That sum is a bound and
not a measure: what a community can carry is what its underwriters can pay from outside
the set being measured, and the two diverge sharply once members declare against standing
the community itself conferred (\S\ref{sec:standing}). Nor can the community over-issue
by increments: every unit of insured credit is backed by a reservation on a specific chain
of stakes running to an underwriter, and when those are exhausted the next acceptance is
uninsured --- visibly, at the point of decision, to the creditor who would bear it. No
aggregate silently inflates.

\subsection{No separate macroprudential brake}

This design has no brake: no system-wide multiplier that throttles new credit when
measured exposure rises. A brake is an exception to the credit rule, applied on a
different measurement on a different schedule --- a second answer to "how much may this
member owe", the question this design exists to have exactly one answer to --- and as a
global multiplier it reprices every member's capacity for conditions none of them created,
squeezing hardest the members with the least to do with the exposure that triggered it.
What a brake was for is achieved structurally: exposure cannot rise without reservations
rising, reservations cannot rise beyond the cut, and stakes decay continuously unless
renewed, so a community whose trade slows sees its aggregate capacity fall without anybody
measuring anything.

\subsection{Where a loss lands}
\label{sec:loss-lands}

A default on an insured obligation is absorbed by the underwriters whose supply carried
it, and by nobody else (\S\ref{sec:substitution}). The loss does \emph{not} travel along
the chain of stakes the reservation ran through: intermediate stakes on a reserved path
are conduits of evidence, not of liability. A stake is placed mechanically at settlement
--- a member who honours a debt thereby records that their creditor stood behind them ---
so liability along the path would be an externality nobody signed, attached to the act of
trading honestly, chilling exactly the trade that builds the graph and most for the members
with the most trade behind them. The sanction therefore lives where the consent lives: at
the source arc, whose owner declared a supply and thereby consented to inherit debts up to
it (Property~\ref{prop:uw-consent}). A member in the middle of a path loses nothing when a
debtor two hops away fails, and owes nothing on their account.

\subsection{What the cascade couples, and what it does not}

The support cascade (\S\ref{sec:medium}) is the one mechanism that couples members'
fortunes, and \textbf{it couples production, not failure}. When a supporter sells, the
proceeds discharge a listed beneficiary's obligations rather than the supporter's own.
What the supporter gives up is that clearing --- the debts of theirs that would otherwise
have closed, or the claim on the buyer they would otherwise have held --- and nothing
attaches to them afterwards. If the beneficiary later fails, the loss falls where
\S\ref{sec:loss-lands} says every loss falls: on the underwriters whose supply carried the
obligation, or on the creditor alone if it was uninsured. \textbf{No mechanism in this
ledger makes a beneficiary's failure a supporter's}, and there is no record for one to be
built on: the cascade writes no attribution at all (\S\ref{sec:medium}).

Three bounds sit on the coupling that does exist. The recursion is capped at depth
$\kConstMaxCascadeDepth$. Each edge is capped at $\kConstNuDrain$ times what that pair has
staked in one another --- zero for a pair that has never settled anything, so a stranger
who lists you drains nothing whether you approved them or not. And both declarations are
required (Property~\ref{prop:benef}), so no member is drawn in without having said so.
The residual exposure is therefore explicit and belongs to the \emph{supporter}: the set
of beneficiaries they have \emph{listed}, visible in their own client, revocable by them,
and bounded above by what they choose to sell, while what the beneficiary gives up by
approving is standing, not money (Property~\ref{prop:benef}). This is a smaller guarantee
than a system-wide containment bound: coupling that members opt into cannot be bounded by
the protocol without overriding the choice they made.

\subsection{Decay as the withdrawal mechanism}

Decay (\S\ref{sec:standing}) is how a community withdraws backing without any member
having to act: standing given once would otherwise be given for ever, and a relationship
nobody currently has would go on being underwritten. Its floor at the live reservation is
what stops the same mechanism dissolving collateral out from under a debt still owed.
